\setlength{\baselineskip}{20pt}
\begin{englishabstract}
Denoising models, as the core mechanism of diffusion frameworks, have shown strong capability in generative modeling, yet their value in discriminative learning remains insufficiently explored. Discriminative tasks fundamentally rely on representation learning, i.e., learning features that support accurate prediction and decision-making. In real visual scenarios, however, noise contamination and distribution perturbation often degrade feature discriminability. Existing diffusion-based discriminative approaches are usually task-specific and often require explicit multi-step denoising during inference, making it difficult to jointly satisfy performance gains and deployment efficiency. To address this gap, this thesis studies discriminative feature enhancement from two perspectives, and the main contributions are summarized as follows:
\par
(1) Feature enhancement without extra inference cost: cascaded embedding layers in a discriminative backbone are reinterpreted as progressive denoising stages, enabling unified modeling of feature extraction and denoising. A parameter fusion mechanism is further derived with strict equivalence between pre-fusion and post-fusion forms, so that training-time dual-branch enhancement can be folded into deployment-time single-branch computation. This design delivers computation-free enhancement without changing task heads or model scale, and is compatible with both label-free and label-augmented training for denoising--discriminative co-optimization. Moreover, it integrates into diverse discriminative backbones with low migration cost and a simple deployment workflow, and can be directly plugged into existing industrial inference pipelines.
\par
(2) Feature enhancement via model distillation: to mitigate the limited representational capacity of lightweight denoising modules, a teacher--student distillation scheme is introduced to provide more reliable noise-direction supervision. In addition, multi-round single-layer denoising and multi-step denoising fusion are proposed to compress multi-step trajectories and strengthen denoising depth without increasing deployment complexity. The two perspectives form a progressive framework in which the first addresses deployment efficiency, while the second further improves robustness and generalization under the same computational budget. Furthermore, distillation and multi-step fusion are jointly optimized in one training pipeline for stable convergence while maintaining controllable optimization behavior.
\par
Experiments on person re-identification, vehicle re-identification, image classification, object detection, and semantic segmentation, covering both CNN and Transformer backbones on multiple public datasets, with metrics including mAP/Rank-1, Top-1/Top-5, AP series, and mIoU, demonstrate consistent gains and favorable efficiency. These results verify strong cross-task generalization and practical deployability, and provide a unified, effective, and actionable pathway for applying diffusion principles to discriminative representation learning, with promising practical and engineering applicability and sustained room for evolution.
\par\noindent\englishkeywords{diffusion model; discriminative representation learning; feature enhancement; parameter fusion; model distillation}
\end{englishabstract}
